\section{项目背景}

\subsection{主题}

本项目以中国古典诗词为素材，覆盖\textbf{公元 712 年（盛唐）至 1279 年（南宋末）}，跨越唐、五代、两宋，
构建一个叙事型数据可视化系统，用可量化的文本证据考察文学如何回应时代变迁。

\subsection{研究立意：探索性，而非预设结论}

我们\textbf{不预设}“唐外向、宋内向”的单一结论，而是以两次最剧烈的历史断裂为锚，做探索性分析：

\begin{keybox}[统领问题]
中国文人在两次大断裂——\textbf{安史之乱（755）}、\textbf{靖康之变（1127）}——前后，
精神世界与文学表达发生了怎样的变化？
\end{keybox}

数据已显示这并非单调的“由外向内”：唐代外向性确在下降，\textbf{但靖康之变后，南宋豪放派
（辛弃疾、陆游的家国激愤）使外向性回潮}，与姜夔、吴文英的醇雅并存——一个断裂岔出两条路。
这种“反弹”正是结论的一部分，而非噪音。

\subsection{为什么是 712–1279 这个跨度}

\begin{itemize}[leftmargin=*, itemsep=0.3em]
    \item \textbf{覆盖完整周期}：兴起—极盛—断裂—沉潜—重构—再断裂—再分化，两次“盛极而崩”互为镜像。
    \item \textbf{两大断裂可对照}：755 与 1127 都是“军事崩溃 + 被迫南迁 + 文学剧变”，并置可见文学回应历史的共性与差异。
    \item \textbf{唐诗 vs 宋词}：完成古典诗词最经典的纵向对比。
    \item \textbf{可验证}：每个转变都有可量化的文本证据（情感、意象、地理、词频多维同步）。
\end{itemize}

\subsection{历史分期：五期 + 两道断裂带}

时间轴为有序时期轴，两道断裂带分别落在 755（盛唐\,$\vert$\,中唐）与 1127（北宋\,$\vert$\,南宋）。

\begin{table}[H]
\centering
\renewcommand{\arraystretch}{1.3}
\begin{tabular}{>{\bfseries}l l l l}
\toprule
时期 & 大致时间 & 社会状态 & 文学精神 \\
\midrule
盛唐 & 712--755 & 开元盛世 & 昂扬、开阔、功业向往 \\
中唐 & 756--835 & 安史之乱后、藩镇 & 沉郁、写实、转向内省 \\
晚唐五代 & 836--960 & 衰乱、分裂 & 感伤、绮丽、向内 \\
北宋 & 960--1127 & 统一、儒学复兴 & 理性、清雅、日常审美（内向最深） \\
南宋 & 1127--1279 & 偏安、抗金 & 豪放回潮 + 醇雅并存 \\
\bottomrule
\end{tabular}
\end{table}

\noindent\textbf{两个人形节点}：杜甫横跨安史之乱（盛唐\,$\to$\,中唐），李清照横跨靖康之变（北宋\,$\to$\,南宋）——
分别是两次断裂亲历的“最亮的星”，将在可视化中作为接缝处的关键个案。
